\section*{摘\quad 要}

本文研究含光伏、储能和外部电网的小区微网购电问题。难点在于光伏与负载随时间变化，电池存在容量、功率和充放电损耗，而普通购电、改单和紧急购电又采用不同的价格规则。为使四问前后连贯，本文先建立统一的十分钟电量流和储能递推关系，再根据各问的信息条件逐步加入供需预测、风险准备、日内预报和波动电价。

针对问题一，建立\textbf{全时段电量流线性规划模型}，统筹光伏直供、外网购电和储能移峰，并用日初、日末储电量相等保证费用比较公平。所得代表日费用为35126.85元，低于峰谷购电策略的38276.64元和光伏空间预留策略的37944.35元。

针对问题二，建立\textbf{周期预测--近期加权准备--跨日线性规划模型}。负载与光伏分别依据历史规律预测，再用近期净需求误差确定安全准备量，以应对五倍价格的紧急购电。2025年2月至12月总费用为13583810.84元，其中普通电费12979560.57元、紧急电费604250.27元；相较等权准备方案节省44.66万元。参数检验表明，0.7的日期衰减在快速适应近期变化和保留足够历史信息之间取得了较好平衡。

针对问题三，建立\textbf{历史融合预报--动态安全准备--日内滚动规划模型}。模型将官方光伏预报与近期历史按过往误差自动融合，并在0:00、6:00、12:00和18:00滚动调整尚未执行的订单。全年总费用为13146535.49元，比不加准备量的初始融合方案节省94.22万元。在日内反馈和改单权限相同的条件下，引入6:00、12:00和18:00的新预报又节省21.78万元，说明后发预报值得使用，但每次是否改单仍需结合调整费和储电状态判断。

针对问题四，将未知波动电价纳入预测和调度。零点一次订电分支采用\textbf{滚动提升树--零点线性规划模型}，总费用为14248570.87元；日内调整分支采用\textbf{滚动岭回归--日内调整线性规划模型}，总费用为13822901.35元。后者虽未取得候选中的绝对最低费用，却以0.016\%的费用差换得更少的紧急购电和更清楚的价格修正依据。逐段回算表明，各方案满足供需平衡、储能容量和充放电功率限制。模型能够在费用、供电风险和解释性之间形成清楚的取舍，可为社区、校园和园区微网的滚动购电提供参考。

\noindent\textbf{关键词：}微网调度；储能优化；滚动规划；风险准备；波动电价
